\documentclass[11pt]{article}
\usepackage[a4paper,margin=28mm]{geometry}
\usepackage{amsmath,amssymb,amsthm}
\usepackage[colorlinks=true,linkcolor=blue,urlcolor=blue,citecolor=blue]{hyperref}
\newtheorem{theorem}{Theorem}
\newcommand{\Q}{\mathbb Q}
\newcommand{\Z}{\mathbb Z}
\newcommand{\F}{\mathbb F}
\DeclareMathOperator{\End}{End}
\DeclareMathOperator{\Gal}{Gal}
\DeclareMathOperator{\GL}{GL}
\DeclareMathOperator{\PGL}{PGL}
\title{No rational \(31\)-torsion on geometrically split abelian surfaces}
\author{}
\date{20 October 2026}

\begin{document}
\maketitle

\begin{theorem}\label{thm:main}
Let \(A/\Q\) be a geometrically split abelian surface. Then
\[
 A(\Q)[31]=0.
\]
In particular, \([31]\) does not occur as the rational torsion subgroup
of a geometrically split genus-\(2\) Jacobian over \(\Q\).
\end{theorem}

Here ``geometrically split'' means that \(A_{\overline{\Q}}\) is
isogenous to a product of elliptic curves. We impose no restriction on
the degree of this isogeny. The theorem does not require a principal
polarization.

\section{Inputs on torsion}

For a positive integer \(d\), let \(S(d)\) be the set of primes that
occur as the order of a torsion point on an elliptic curve over a
number field of degree \(d\). We use the following results.
\begin{enumerate}
\item Derickx--Kamienny--Stein--Stoll \cite[Theorem 1.2]{DKSS} prove that
\[
 S(4)=\{2,3,5,7,11,13,17\},\qquad
 S(6)=\{2,3,5,7,11,13,17,19,37\}.
\]
Hence an elliptic curve over a field of degree \(1,2,3,4,\) or \(6\)
has no point of order \(31\). For degrees \(1\) and \(2\), extend
to a field of degree \(4\); for degree \(3\), extend to one of degree
\(6\).
\item Khawaja \cite[Theorem 1.2]{Khawaja} proves that
\[
 S(8)=\{2,3,5,7,11,13,17,19,23\}.
\]
Thus an elliptic curve over a degree-\(8\) number field has no point
of order \(31\).
\item Clark \cite[Theorem 5]{Clark} proves that, if an abelian surface
over \(\Q\) has everywhere potentially good reduction, its rational
torsion order belongs to
\[
 \{1,\ldots,16,18,19,20,22,24,25,28,30,36,48,60,72\}.
\]
None of these integers is divisible by \(31\).
\end{enumerate}

We also use the standard independence of the semisimplified residual
representation from the choice of a stable lattice. It applies to
isomorphic \(p\)-adic representations after transporting a lattice
through the isomorphism.

\section{The elementary cases and the endomorphism field}

Suppose that \(P\in A(\Q)\) has order \(31\).
First suppose that \(A_{\overline{\Q}}\) is isogenous to
\(E_1\times E_2\), where the elliptic curves \(E_1,E_2\) are not
geometrically isogenous. There are exactly two geometric elliptic
subvarieties of \(A\). Galois permutes them, so both are defined over
a field \(K\) with \([K:\Q]\leq2\). Choose one of them, \(E\),
and put \(F=A_K/E\). The exact sequence
\[
 0\longrightarrow E\longrightarrow A_K\longrightarrow F
 \longrightarrow0
\]
shows that \(P\) either belongs to \(E(K)\) or maps to a point of
order \(31\) in \(F(K)\). Both possibilities contradict the elliptic
torsion bounds above.

Next suppose that \(A_{\overline{\Q}}\) is isogenous to \(E^2\),
where \(E\) has complex multiplication. CM elliptic curves have
everywhere potentially good reduction. This property is preserved by
products and isogenies, and descends from a finite extension of the
ground field. Hence \(A\) has everywhere potentially good reduction,
and Clark's theorem gives a contradiction.

It remains to consider the case where \(A_{\overline{\Q}}\) is
isogenous to \(E^2\) for a non-CM elliptic curve \(E\). Then
\[
 D=\End^0(A_{\overline{\Q}})\simeq M_2(\Q).
\]
Let \(L/\Q\) be the smallest field over which all geometric
endomorphisms of \(A\) are defined. It is a finite Galois extension.
Put \(\Gamma=\Gal(L/\Q)\). By Skolem--Noether, its faithful action
on \(D\) embeds \(\Gamma\) into \(\PGL_2(\Q)\).

A finite subgroup of \(\PGL_2(\mathbb R)\) is cyclic or dihedral:
it fixes a point of the hyperbolic plane, and its action on the tangent
plane embeds it in \(O(2)\). If an element of \(\PGL_2(\Q)\) has
order \(n>2\), the ratio of the eigenvalues of a representative \(M\)
is a primitive \(n\)-th root of unity \(\zeta\). Thus
\[
 \frac{(\operatorname{tr}M)^2}{\det M}
 =2+\zeta+\zeta^{-1}\in\Q,
\]
which implies \(n\in\{3,4,6\}\). Hence \(\Gamma\) is cyclic of
order \(1,2,3,4,\) or \(6\), or dihedral of order \(4,6,8,\) or
\(12\). Here the dihedral group of order \(4\) is the Klein four
group. In particular,
\[
 [L:\Q]\in\{1,2,3,4,6,8,12\}.
\]

Every geometric elliptic subvariety of \(A\) is defined over \(L\):
it is the image of an integral multiple of an idempotent in \(D\).
Applying the elliptic exact sequence over \(L\) and the torsion
bounds from Section~1 excludes every listed degree except \(12\).
We are therefore left with
\[
 \Gamma=D_{12},
\]
where the subscript denotes the group order. Its rotation subgroup
\(C=\langle r\rangle\) has order \(6\). Put \(K=L^C\), so
\([K:\Q]=2\).

\section{A stable split quadratic order}

Set \(p=31\), \(k=\Q_p\), and \(\mathcal O=\Z_p\). Identify
\(D\otimes_{\Q}k\) with \(M_2(k)\), and normalize a matrix
representative \(R\) of \(r\) to have trace \(1\). Since \(r\)
has projective order \(6\), the characteristic polynomial of \(R\) is
\[
 X^2-X+\frac13.
\]
Its discriminant \(-1/3\) is a nonzero square modulo \(31\): it is
\(10\) modulo \(31\), and \(14^2\equiv10\pmod{31}\). Thus its
eigenvalues \(\alpha,\beta\) lie in \(\mathcal O^\times\), and
\(\alpha-\beta\) is a unit. It follows that
\[
 \mathcal B=\mathcal O[R]\simeq\mathcal O\times\mathcal O.
\]
Write \(e_1,e_2\) for its two orthogonal idempotents.

Conjugation by \(\Gamma\) sends \(r\) to \(r\) or \(r^{-1}\).
The trace-normalized representative of \(r^{-1}\) is \(1-R\), so
\(\Gamma\) preserves \(\mathcal B\) and permutes
\(\{e_1,e_2\}\). The rotation subgroup fixes each idempotent, and
every element outside it exchanges them.

Let \(T=T_p(A)\), and set
\[
 T'=\mathcal B T.
\]
Since \(\mathcal B\) is finite free over \(\mathcal O\), this is a
full lattice in \(V_p(A)\). It is stable under
\(G_{\Q}=\Gal(\overline{\Q}/\Q)\), since both \(T\) and
\(\mathcal B\) are stable. We obtain
\[
 T'=T_1\oplus T_2,\qquad T_i=e_iT'.
\]
Each \(T_i\) is free of rank \(2\) over \(\mathcal O\): the
four-dimensional \(M_2(k)\)-module \(V_p(A)\) is the sum of two
copies of its standard module. The group \(G_K\) preserves each
\(T_i\), while every element of \(G_{\Q}\setminus G_K\) exchanges
the two summands.

\section{The two summands differ by a character}

Put \(V_i=T_i\otimes_{\mathcal O}k\), and choose an off-diagonal
matrix unit \(X\in e_1M_2(k)e_2\). The action of \(G_K\) on
\(D\otimes k\) is through \(C\). Consequently
\[
 g(X)=\eta(g)X\qquad(g\in G_K),
\]
where \(\eta:G_K\to k^\times\) is the eigenvalue-ratio character,
with value \(\alpha/\beta\) on a lift of \(r\). It factors through
\(\Gal(L/K)=C\) and has exact order \(6\). Its values are units,
and its reduction \(\overline\eta\) also has exact order \(6\),
since reduction is injective on roots of unity of order prime to \(31\).

The matrix unit \(X:V_2\to V_1\) is an isomorphism, and the equality
\(gXg^{-1}=\eta(g)X\) gives
\[
 V_1\simeq\eta\otimes V_2
\]
as \(G_K\)-representations. Put \(M_i=T_i/pT_i\). Applying
lattice independence to this isomorphism gives
\[
 M_1^{\mathrm{ss}}\simeq
 \overline\eta\otimes M_2^{\mathrm{ss}}.
\]
This step does not require \(X\) to carry \(T_2\) onto \(T_1\)
integrally.

\section{The determinant contradiction}

The point \(P\) generates a trivial submodule of
\(A[p]=T/pT\). By lattice independence, \(T'/pT'\) therefore has a
trivial \(G_{\Q}\)-composition factor. Restriction to \(G_K\) is
exact, so at least one of \(M_1,M_2\) has a trivial
\(G_K\)-composition factor. An element of \(G_{\Q}\setminus G_K\)
conjugates these two \(G_K\)-modules. Since a trivial module remains
trivial under conjugation, both have a trivial composition factor.

The relation
\(M_1^{\mathrm{ss}}\simeq\overline\eta\otimes M_2^{\mathrm{ss}}\)
shows that \(\overline\eta^{-1}\) is also a composition factor of
\(M_2\). Since \(M_2\) has dimension \(2\) and
\(\overline\eta\) is nontrivial, we obtain
\[
 M_2^{\mathrm{ss}}\simeq\mathbf1\oplus\overline\eta^{-1},
 \qquad
 M_1^{\mathrm{ss}}\simeq\overline\eta\oplus\mathbf1.
\]
Hence
\[
 \det(T'/pT')|_{G_K}=\det(M_1)\det(M_2)=1.
\]

Every abelian variety over \(\Q\) has a polarization over \(\Q\).
The induced nondegenerate alternating pairing on \(V_p(A)\) shows
that the determinant of its \(p\)-adic representation is the square
of the \(p\)-adic cyclotomic character. A change of stable lattice
does not change this determinant. Writing
\(\chi_p:G_{\Q}\to\F_p^\times\) for the mod-\(p\) cyclotomic
character, we therefore have
\[
 \chi_{31}^2|_{G_K}=1.
\]
The character \(\chi_{31}^2\) has order \(15\), and its kernel
fixes \(\Q(\zeta_{31}+\zeta_{31}^{-1})\), a field of degree \(15\)
over \(\Q\). The displayed identity would put this field inside the
quadratic field \(K\), a contradiction. This proves
Theorem~\ref{thm:main}.
\hfill\(\square\)

\section{Consequence for the census}

The group
\href{https://genus-2-torsion.github.io/group.html?g=31}{\([31]\)}
is impossible in both the class ``split over \(\Q\)'' and the class
``geometrically split, simple over \(\Q\)''.
The argument applies to every geometric splitting degree, including
degrees divisible by \(31\). It uses a stable split quadratic order
at \(31\), without requiring a stable maximal order or a maximal
integral endomorphism ring.

\begin{thebibliography}{9}
\bibitem{DKSS}
M. Derickx, S. Kamienny, W. Stein and M. Stoll,
\emph{Torsion points on elliptic curves over number fields of small
degree}, Algebra \& Number Theory \textbf{17} (2023), 267--308, Theorem 1.2.
\href{https://arxiv.org/pdf/1707.00364}{arXiv:1707.00364}.

\bibitem{Khawaja}
M. Khawaja,
\emph{Torsion primes for elliptic curves over degree \(8\) number fields},
Research in Number Theory \textbf{10} (2024), article 48, Theorem 1.2.
\href{https://eprints.whiterose.ac.uk/id/eprint/211212/8/s40993-024-00533-6.pdf}{Published article}.

\bibitem{Clark}
P. L. Clark,
\emph{Bounds for torsion on abelian varieties with integral moduli},
arXiv:math/0407264v2 (2004), Theorem 5.
\href{https://arxiv.org/pdf/math/0407264}{arXiv:math/0407264}.
The same list is Theorem~1.4 of P. L. Clark and X. Xarles, \emph{Local bounds for torsion points on abelian varieties}, Canad. J. Math. \textbf{60} (2008), 532--555, \href{https://doi.org/10.4153/CJM-2008-026-x}{doi:10.4153/CJM-2008-026-x}, for abelian surfaces with everywhere anisotropic reduction, which includes potentially good reduction.
\end{thebibliography}

\end{document}
